\documentclass{ximera}
%verbeterd 24 september 2026
\title{Vraagstukken elektrisch veld: basis}
\author{Daan Lipkens}

\begin{document}
\begin{abstract}
    Basisoefeningen op het elektrisch veld van een puntlading.
\end{abstract}
\maketitle
%———————————————————————————————————————
% OEF E=F/Q 1: F berekenen uit Q en E (nC en kN/C)
\begin{exercise}
    Een testlading $Q_t = +50 \, \mathrm{nC}$ bevindt zich in een elektrisch veld van $E = 2{,}0 \, \mathrm{kN/C}$.

    Bereken de grootte van de elektrische kracht op de testlading.

    \[ F_e = \answer[tolerance=0.005,onlineshowanswerbutton]{1.0}\cdot 10^{-4} \, \mathrm{N} \]

    \begin{solution}\nl
        \underline{Gegevens:}
        \begin{align*}
            & Q_t = +50 \, \mathrm{nC} = 5{,}0 \cdot 10^{-8} \, \mathrm{C} & & E = 2{,}0 \, \mathrm{kN/C} = 2{,}0 \cdot 10^{3} \, \mathrm{N/C}
        \end{align*}

        \underline{Gevraagd:} $F_e$\\

        \underline{Oplossing:}\\
        \begin{align*}
            E &= \frac{F_e}{Q_t} \\
            \Leftrightarrow F_e &= Q_t \cdot E \\
                &= 5{,}0 \cdot 10^{-8} \, \mathrm{C} \cdot 2{,}0 \cdot 10^{3} \, \mathrm{N/C} \\
                &= 1{,}0 \cdot 10^{-4} \, \mathrm{N} 
        \end{align*}
    \end{solution}
\end{exercise}
%———————————————————————————————————————

%———————————————————————————————————————
% OEF E=F/Q 2: E berekenen uit F en Q (mN en nC)
\begin{exercise}
    Een testlading $Q_t = +6{,}0 \, \mathrm{nC}$ ondervindt in een bepaald punt een elektrische kracht van $3{,}0 \, \mathrm{mN}$.

    Bereken de grootte van het elektrisch veld in dat punt.

    \[ E = \answer[tolerance=0.05,onlineshowanswerbutton]{5.0} \cdot 10^{5} \, \mathrm{N/C} \]

    \begin{solution}\nl
        \underline{Gegevens:}
        \begin{align*}
            & Q_t = +6{,}0 \, \mathrm{nC} = 6{,}0 \cdot 10^{-9} \, \mathrm{C} & & F_e = 3{,}0 \, \mathrm{mN} = 3{,}0 \cdot 10^{-3} \, \mathrm{N}
        \end{align*}

        \underline{Gevraagd:} $E$\\

        \underline{Oplossing:}\\
        \begin{align*}
            E &= \frac{F_e}{Q_t} \\
              &= \frac{3{,}0 \cdot 10^{-3} \, \mathrm{N}}{6{,}0 \cdot 10^{-9} \, \mathrm{C}} \\
              &= 5{,}0 \cdot 10^{5} \, \mathrm{N/C}
        \end{align*}
    \end{solution}
\end{exercise}
%———————————————————————————————————————

%———————————————————————————————————————
% OEF E=F/Q 3: Q berekenen uit F en E (mN en kN/C)
\begin{exercise}
    Op een onbekende testlading werkt in een elektrisch veld van $E = 150 \, \mathrm{kN/C}$ een elektrische kracht van $45 \, \mathrm{mN}$.

    Bereken de grootte van de testlading $Q_t$.

    \[ Q_t = \answer[tolerance=0.01,onlineshowanswerbutton]{3.0}\cdot 10^{-7} \, \mathrm{\mu C} \]

    \begin{hint}
        Herschrijf $E = F_e/Q_t$ naar $Q_t$.
    \end{hint}

    \begin{solution}\nl
        \underline{Gegevens:}
        \begin{align*}
            & F_e = 45 \, \mathrm{mN} = 4{,}5 \cdot 10^{-2} \, \mathrm{N} & & E = 150 \, \mathrm{kN/C} = 1{,}5 \cdot 10^{5} \, \mathrm{N/C}
        \end{align*}

        \underline{Gevraagd:} $Q_t$\\

        \underline{Oplossing:}\\
        \begin{align*}
            E &= \frac{F_e}{Q_t} \\
            \Leftrightarrow Q_t &= \frac{F_e}{E} \\
                &= \frac{4{,}5 \cdot 10^{-2} \, \mathrm{N}}{1{,}5 \cdot 10^{5} \, \mathrm{N/C}} \\
                &= 3{,}0 \cdot 10^{-7} \, \mathrm{C}
        \end{align*}
    \end{solution}
\end{exercise}
%———————————————————————————————————————

%———————————————————————————————————————
% OEF E=F/Q 4: negatieve lading, richting van de kracht
\begin{exercise}
    Een lading $Q = -4{,}0 \, \mathrm{\mu C}$ bevindt zich in een elektrisch veld $\vec{E}$ dat naar rechts wijst, met grootte $E = 2{,}5 \cdot 10^{4} \, \mathrm{N/C}$.

    \begin{question}
        Bereken de grootte van de elektrische kracht op de lading.

        \[ F_e = \answer[tolerance=0.01,onlineshowanswerbutton]{0.10} \, \mathrm{N} \]

        \begin{solution}\nl
            \underline{Gegevens:}
            \begin{align*}
                & Q = -4{,}0 \, \mathrm{\mu C} = -4{,}0 \cdot 10^{-6} \, \mathrm{C} & & E = 2{,}5 \cdot 10^{4} \, \mathrm{N/C}
        \end{align*}

        \underline{Gevraagd:} $F_e$\\

        \underline{Oplossing:}\\            
            \begin{align*}
             E &= \frac{F_e}{Q_t} \\
            \Leftrightarrow F_e &= |Q| \cdot E \\
            &= 4{,}0 \cdot 10^{-6} \, \mathrm{C} \cdot 2{,}5 \cdot 10^{4} \, \mathrm{N/C} \\
            &= 0{,}10 \, \mathrm{N}\\
            & = 1{,}0\cdot 10^{-1} \, \mathrm{N}
            \end{align*}
        \end{solution}
    \end{question}

    \begin{question}
        Naar waar wijst deze kracht?

        \wordChoice{\choice{Naar rechts, dezelfde zin als $\vec{E}$}\choice[correct]{Naar links, tegengesteld aan $\vec{E}$}}

        \begin{hint}
            Vergelijk met het teken van de lading: wijzen kracht en veld voor een negatieve lading in dezelfde of in tegengestelde zin?
        \end{hint}

        \begin{solution}\nl
            Omdat de lading \textbf{negatief} is, wijst de kracht $\vec{F}_e$ altijd in de \textbf{tegengestelde} zin van het veld $\vec{E}$. Het veld wijst naar rechts, dus de kracht wijst naar \textbf{links}.
        \end{solution}
    \end{question}
\end{exercise}
%———————————————————————————————————————

%———————————————————————————————————————
% OEF E=F/Q 5: teken van de lading afleiden uit de waargenomen kracht
\begin{exercise}
    In een bepaald punt heerst een elektrisch veld $\vec{E}$ dat \textbf{naar boven} wijst, met grootte $E = 8{,}0 \cdot 10^{3} \, \mathrm{N/C}$. Een lading die in dat punt wordt geplaatst, ondervindt een kracht van $16 \, \mathrm{mN}$, gericht \textbf{naar onder}.

    \begin{question}
        Welk teken heeft deze lading?

        \wordChoice{\choice{Positief}\choice[correct]{Negatief}}

        \begin{hint}
            Vergelijk de zin van de kracht met de zin van het veld.
        \end{hint}

        \begin{solution}\nl
            De kracht wijst naar onder, terwijl het veld naar boven wijst: kracht en veld hebben dus een \textbf{tegengestelde} zin. Dat is enkel mogelijk bij een \textbf{negatieve} lading.
        \end{solution}
    \end{question}

    \begin{question}
        Bereken de grootte van de lading.

        \[ |Q| = \answer[tolerance=0.1,onlineshowanswerbutton]{2.0}\cdot 10^{-6} \, \mathrm{C} \]

        \begin{solution}\nl
            \underline{Gegevens:}
            \begin{align*}
                & F_e = 16 \, \mathrm{mN} = 1{,}6 \cdot 10^{-2} \, \mathrm{N} & & E = 8{,}0 \cdot 10^{3} \, \mathrm{N/C}
            \end{align*}

            \underline{Gevraagd:} $|Q|$\\

            \underline{Oplossing:}\\
            \begin{align*}
             E &= \frac{F_e}{Q_t} \\
            \Leftrightarrow |Q| &= \frac{F_e}{E} \\
                    &= \frac{1{,}6 \cdot 10^{-2} \, \mathrm{N}}{8{,}0 \cdot 10^{3} \, \mathrm{N/C}} \\
                    &= 2{,}0 \cdot 10^{-6} \, \mathrm{C}
            \end{align*}
        \end{solution}
    \end{question}
\end{exercise}
%———————————————————————————————————————

%———————————————————————————————————————
% OEF BASIS 6: E van puntlading
\begin{exercise}
    Bereken de grootte van het elektrisch veld op een afstand van $0{,}50 \, \mathrm{m}$ van een puntlading $Q = +6{,}0 \cdot 10^{-6} \, \mathrm{C}$.

    \[ E = \answer[tolerance=0.05,onlineshowanswerbutton]{2.2} \cdot 10^{5} \, \mathrm{N/C} \]

    \begin{solution}\nl
        \underline{Gegevens:}
        \begin{align*}
            & Q = +6{,}0 \cdot 10^{-6} \, \mathrm{C} & & r = 0{,}50 \, \mathrm{m}
        \end{align*}

        \underline{Gevraagd:} $E$\\

        \underline{Oplossing:}\\
        \begin{align*}
            E &= k_e \frac{|Q|}{r^{2}} \\
              &= \frac{8{,}99 \cdot 10^{9} \, \mathrm{\frac{N \cdot m^2}{C^2}} \cdot 6{,}0 \cdot 10^{-6} \, \mathrm{C}}{(0{,}50 \, \mathrm{m})^{2}} \\
              &= 2{,}2 \cdot 10^{5} \, \mathrm{N/C}
        \end{align*}
    \end{solution}
\end{exercise}
%———————————————————————————————————————

%———————————————————————————————————————
% OEF BASIS 7: E van negatieve puntlading + richting
\begin{exercise}
    Een puntlading $Q = -3{,}0 \cdot 10^{-6} \, \mathrm{C}$ bevindt zich op een afstand van $0{,}20 \, \mathrm{m}$ van punt \textit{P}.

    \begin{question}
        Bereken de grootte van het elektrisch veld $E$ in punt \textit{P}.
        \[ E = \answer[tolerance=0.05,onlineshowanswerbutton]{6.7} \cdot 10^{5} \, \mathrm{N/C} \]

        \begin{solution}\nl
            \[
                E = k_e \frac{|Q|}{r^{2}} = \frac{8{,}99 \cdot 10^{9} \, \mathrm{\frac{N \cdot m^2}{C^2}} \cdot 3{,}0 \cdot 10^{-6} \, \mathrm{C}}{(0{,}20 \, \mathrm{m})^{2}} = 6{,}7 \cdot 10^{5} \, \mathrm{N/C}
            \]
        \end{solution}
    \end{question}

    \begin{question}
        Wat is de richting en zin van het elektrisch veld in punt \textit{P}?

        \wordChoice{\choice{Van de lading weg}\choice[correct]{Naar de lading toe}}
        \begin{onlineOnly}
            \begin{hint}
                Wat is het teken van de bronlading?
            \end{hint}
        \end{onlineOnly}
        \begin{solution}\nl
            De bronlading is \textbf{negatief}, dus wijst het elektrisch veld radiaal \textbf{naar de lading toe}.
        \end{solution}
    \end{question}
\end{exercise}
%———————————————————————————————————————

%———————————————————————————————————————
% OEF BASIS 8: E uit F en Q
\begin{exercise}
    Een testlading $Q_t = +2{,}0 \cdot 10^{-6} \, \mathrm{C}$ ondervindt in een bepaald punt een elektrische kracht van $0{,}015 \, \mathrm{N}$. Bereken de grootte van het elektrisch veld in dat punt.

    \[ E = \answer[tolerance=100,onlineshowanswerbutton]{7500} \, \mathrm{N/C} \]


    \begin{solution}\nl
        \underline{Gegevens:}
        \begin{align*}
            & Q_t = +2{,}0 \cdot 10^{-6} \, \mathrm{C} & & F_e = 0{,}015 \, \mathrm{N}
        \end{align*}

        \underline{Gevraagd:} $E$\\

        \underline{Oplossing:}\\
        \begin{align*}
            E &= \frac{F_e}{Q_t} \\
              &= \frac{0{,}015 \, \mathrm{N}}{2{,}0 \cdot 10^{-6} \, \mathrm{C}} \\
              &= 7{,}5 \cdot 10^{3} \, \mathrm{N/C}
        \end{align*}
    \end{solution}
\end{exercise}
%———————————————————————————————————————

%———————————————————————————————————————
% OEF BASIS 9: F uit E en Q
\begin{exercise}
    In een bepaald punt heerst een elektrisch veld van $E = 4{,}0 \cdot 10^{4} \, \mathrm{N/C}$. Een lading $Q = +5{,}0 \cdot 10^{-6} \, \mathrm{C}$ wordt in dat punt geplaatst.

    \[ F_e = \answer[tolerance=0.01,onlineshowanswerbutton]{0.20} \, \mathrm{N} \]

    \begin{hint}
        Gebruik $F_e = Q \cdot E$.
    \end{hint}

    \begin{solution}\nl
        \underline{Gegevens:}
        \begin{align*}
            & E = 4{,}0 \cdot 10^{4} \, \mathrm{N/C} & & Q = +5{,}0 \cdot 10^{-6} \, \mathrm{C}
        \end{align*}

        \underline{Gevraagd:} $F_e$\\

        \underline{Oplossing:}\\
        \begin{align*}
             E &= \frac{F_e}{Q_t} \\
            \Leftrightarrow F_e &= Q \cdot E \\
                &= 5{,}0 \cdot 10^{-6} \, \mathrm{C} \cdot 4{,}0 \cdot 10^{4} \, \mathrm{N/C} \\
                &= 0{,}20 \, \mathrm{N}= 2{,}0\cdot 10^{-1} \, \mathrm{N}
        \end{align*}
        Omdat $Q$ positief is, heeft $F_e$ dezelfde richting en zin als $\vec{E}$.
    \end{solution}
\end{exercise}
%———————————————————————————————————————

%———————————————————————————————————————
% OEF BASIS 10: r bepalen uit gegeven E
\begin{exercise}
    Op welke afstand van een puntlading $Q = +8{,}0 \cdot 10^{-6} \, \mathrm{C}$ is de grootte van het elektrisch veld gelijk aan $2{,}0 \cdot 10^{5} \, \mathrm{N/C}$?

    \[ r = \answer[tolerance=0.02,onlineshowanswerbutton]{0.60} \, \mathrm{m} \]

    \begin{hint}
        Herschrijf de formule $E = k_e \cdot |Q|/r^2$ naar $r$.
    \end{hint}

    \begin{solution}\nl
        \underline{Gegevens:}
        \begin{align*}
            & Q = +8{,}0 \cdot 10^{-6} \, \mathrm{C} & & E = 2{,}0 \cdot 10^{5} \, \mathrm{N/C}
        \end{align*}

        \underline{Gevraagd:} $r$\\

        \underline{Oplossing:}\\
        \begin{align*}
            E &= k_e \frac{|Q|}{r^{2}} \\
            r^{2} &= k_e \frac{|Q|}{E} = \frac{8{,}99 \cdot 10^{9} \, \mathrm{\frac{N \cdot m^2}{C^2}} \cdot 8{,}0 \cdot 10^{-6} \, \mathrm{C}}{2{,}0 \cdot 10^{5} \, \mathrm{N/C}} = 0{,}36 \, \mathrm{m^2}\\
            r &= 0{,}60 \, \mathrm{m} = 6{,}0\cdot 10^{-1} \mathrm{m}
        \end{align*}
    \end{solution}
\end{exercise}
%———————————————————————————————————————

\end{document}